\chapter{Introducción, motivación y objetivos}
\label{cap:introduccion}

El presente capítulo introduce el contexto general en el que se enmarca este
Trabajo Fin de Grado, describe la motivación que justifica su desarrollo y
delimita los objetivos que guían la solución propuesta. Para ello, en primer
lugar se analiza el papel del Currículum Vítae Normalizado (CVN) en la gestión de
información académica e investigadora en España. A continuación, se expone la
necesidad de disponer de una representación de bajo nivel que facilite el
procesamiento automático de los datos curriculares. Finalmente, se establecen el
alcance del proyecto, las competencias de Computación trabajadas, las principales
contribuciones realizadas y la estructura seguida por el resto del documento.

\section{Contexto}
\label{sec:introduccion_contexto}

La actividad académica e investigadora produce una cantidad considerable de
información curricular de naturaleza heterogénea. En un mismo currículo pueden
aparecer datos relativos a formación, experiencia profesional, publicaciones,
proyectos de investigación, contratos, estancias, dirección de trabajos,
transferencia, actividad docente y otros méritos que responden a estructuras y
criterios de descripción diferentes. Esta información no se utiliza únicamente
como documento de presentación personal, sino también como entrada para procesos
de evaluación, convocatorias públicas, solicitudes de ayudas, acreditaciones,
portales institucionales y sistemas internos de gestión de la investigación. En
consecuencia, el valor de un currículo no depende solo de su apariencia final,
sino también de la capacidad de los sistemas informáticos para interpretar,
validar, reutilizar y transformar los datos que contiene.

En el ámbito español, el Currículum Vítae Normalizado (CVN) constituye una norma
de referencia para la presentación de información curricular de investigadores. La
Fundación Española para la Ciencia y la Tecnología (FECYT) lo define como una
norma estándar que establece un formato común de presentación de los datos
curriculares y facilita la interoperabilidad con bases de datos
institucionales~\cite{cvn_web}. Esta función normalizadora resulta especialmente
relevante en un entorno en el que una misma trayectoria académica puede tener que
presentarse ante organismos, plataformas y convocatorias con requisitos distintos.
Por ello, CVN representa un punto de partida adecuado para estudiar cómo puede
formalizarse la información curricular de forma computable y reutilizable.

No obstante, la existencia de un formato normalizado no elimina por sí misma el
coste asociado a la elaboración, el mantenimiento y la adaptación de los
currículos. A pesar de disponer de una referencia común, los profesores e
investigadores universitarios continúan dedicando un tiempo significativo a
actualizar méritos, preparar versiones específicas, adaptar contenidos a formatos
concretos o trasladar información entre herramientas. Esta situación se ve
agravada cuando el mismo historial debe reutilizarse en contextos diferentes, ya
que pequeñas variaciones en la estructura, el nivel de detalle o el formato de
salida pueden obligar a realizar tareas manuales repetitivas. Desde esta
perspectiva, resulta necesario avanzar hacia herramientas abiertas que automaticen
parte del proceso y que permitan gestionar los currículos de manera más
sistemática.

El tratamiento computacional de CVN, sin embargo, no puede reducirse a la
generación de un documento final ni a la lectura de un único fichero XML. La base
del trabajo es la norma CVN y el ecosistema de artefactos que la acompaña; dentro
de dicho ecosistema, XML actúa como mecanismo de representación e intercambio de
datos. En particular, el paquete oficial analizado en este trabajo incluye
esquemas XML/XSD, manuales estructurados, modelos auxiliares y tablas de
referencia que describen la estructura de los datos, sus códigos internos y su
organización jerárquica. Estos elementos aportan la base formal necesaria para la
interoperabilidad, pero también introducen una complejidad que dificulta su uso
directo como modelo interno de herramientas abiertas, ligeras y reproducibles. En
conjunto, el problema abordado consiste en separar el significado de la
información curricular de la forma concreta en que la norma la representa y la
serializa en sus artefactos técnicos.

\section{Motivación}
\label{sec:introduccion_motivacion}

La motivación principal de este proyecto surge de la necesidad de definir una
representación de bajo nivel para currículos académicos e investigadores basada en
la norma CVN, pero orientada al procesamiento automático de datos. Dicha
representación debe permitir que la información curricular pueda almacenarse,
validarse, transformarse y exportarse sin depender exclusivamente de la
representación XML empleada por CVN ni de la herramienta oficial asociada. Esta
necesidad está directamente relacionada con la posibilidad de desarrollar
herramientas abiertas para la creación, actualización y adaptación de currículos,
tal y como se plantea en la descripción oficial del Trabajo Fin de Grado.

Antes de entrar en el detalle técnico, conviene explicar esta motivación en
los mismos términos de uso final con los que se planteó originalmente el
proyecto. La idea de partida era avanzar hacia una situación en la que:

\begin{itemize}
  \item \textbf{Existiera una especificación abierta de currículo}, que
  sirviera de referencia común independiente de cualquier herramienta
  concreta.

  \item \textbf{Pudieran existir almacenes de currículos abiertos}
  construidos sobre esa especificación, con una forma bien definida de leer
  y escribir en ellos.

  \item \textbf{Pudieran desarrollarse herramientas de lectura y formato}
  capaces de generar, a partir de un currículo expresado en esa
  especificación, currículos adaptados a necesidades concretas.

  \item \textbf{Pudieran desarrollarse herramientas de importación}
  capaces de incorporar a esa especificación currículos ya existentes,
  creados con la herramienta oficial de FECYT.
\end{itemize}

Los objetivos específicos de la Sección~\ref{sec:introduccion_objetivos} y la
arquitectura Open CVN presentada en el resto del documento traducen estos
cuatro pasos en tareas concretas de Computación, y el
Capítulo~\ref{cap:conclusiones} retoma esta misma motivación para valorar en
qué medida se ha alcanzado.

Desde el punto de vista técnico, el uso directo de los artefactos oficiales de CVN
como modelo interno presenta varias dificultades. Los esquemas XSD proporcionan
una descripción estructural precisa de la representación XML, pero no siempre
expresan de forma explícita la semántica curricular necesaria para construir
aplicaciones de alto nivel. Además, parte de la información relevante se encuentra
distribuida entre manuales, modelos auxiliares y tablas de referencia, lo que
obliga a combinar varias fuentes para interpretar correctamente los campos,
códigos y relaciones del dominio. Esta dispersión no invalida la norma CVN ni sus
artefactos oficiales, que cumplen una función esencial de interoperabilidad, pero
sí condiciona su reutilización como base directa para herramientas de gestión,
validación o exportación.

El presente Trabajo Fin de Grado aborda este problema desde una perspectiva de
Computación. La cuestión central no consiste únicamente en leer o escribir
ficheros, sino en analizar fuentes formales heterogéneas, extraer evidencia
estructural y semántica, normalizar metadatos, preservar trazabilidad y generar
modelos capaces de representar información curricular compleja. De este modo, el
proyecto se sitúa en el ámbito de la formalización y representación computable del
conocimiento, con especial atención a la separación entre formato de
serialización, modelo conceptual y artefactos de validación.

Como respuesta a esta situación, Open CVN propone una base técnica abierta,
extensible y potencialmente colaborativa para representar currículos académicos
tomando CVN como referencia. La solución no pretende sustituir la norma oficial,
sino construir sobre ella una capa intermedia que facilite el desarrollo de
herramientas posteriores. Esta capa debe permitir trabajar con un formato final
basado en JSON, adecuado tanto para almacenamiento en archivos de texto como para
su integración en bases de datos documentales o aplicaciones de procesamiento. De
este modo, la norma CVN y sus artefactos oficiales se mantienen como fuentes de
trazabilidad e interoperabilidad, mientras que el formato Open CVN se orienta a
facilitar la gestión automatizada y la evolución controlada de los datos
curriculares sin quedar limitado a la representación XML ni a la herramienta
oficial existente.

\section{Objetivos}
\label{sec:introduccion_objetivos}

El objetivo general de este Trabajo Fin de Grado es \textit{diseñar e implementar
una arquitectura computacional abierta para representar, validar, transformar,
almacenar y exportar currículos académicos e investigadores tomando CVN como norma
y ecosistema de referencia, pero evitando que su representación XML o la
herramienta oficial asociada condicionen por completo el modelo interno del
sistema, con el fin de facilitar el desarrollo de herramientas abiertas para la
gestión automatizada de currículos}. Para alcanzar esta meta global se han
definido los siguientes objetivos específicos:

\begin{enumerate}
  \item \textbf{Analizar el ecosistema CVN y sus artefactos oficiales.} Este
  objetivo comprende el estudio de los esquemas XML/XSD, manuales estructurados,
  modelos auxiliares y tablas de referencia que forman parte del paquete empleado
  como fuente del proyecto.

  \item \textbf{Estudiar alternativas de representación y validación de datos
  curriculares.} Se consideran modelos conceptuales, formatos de serialización y
  herramientas de generación de modelos con el fin de seleccionar una estrategia
  adecuada para información curricular compleja.

  \item \textbf{Generar una capa estructural reproducible a partir de los XSD
  oficiales.} La generación automática permite obtener artefactos consistentes con
  la representación XML definida dentro del ecosistema CVN y, al mismo tiempo,
  separar el código generado de la lógica mantenida manualmente.

  \item \textbf{Normalizar metadatos procedentes de las fuentes oficiales.} Este
  objetivo busca unificar la información funcional y técnica asociada a los
  códigos CVN, reduciendo la dispersión existente entre documentos auxiliares.

  \item \textbf{Definir reglas semánticas conservadoras y trazables.} Las reglas
  deben permitir interpretar campos, tipos, tablas de referencia y relaciones sin
  perder la conexión con las fuentes originales de las que se derivan.

  \item \textbf{Generar modelos de dominio y artefactos conceptuales.} Estos
  modelos deben representar la información curricular de forma más próxima al
  dominio académico que a la estructura concreta de la serialización XML.

  \item \textbf{Definir un formato Open CVN JSON.} El formato propuesto debe ser
  validable, extensible y adecuado para intercambio, almacenamiento, colaboración
  y transformación posterior.

  \item \textbf{Implementar mecanismos de parseo, validación, almacenamiento y
  exportación.} Estos mecanismos permiten demostrar que la representación definida
  puede emplearse en flujos reales de gestión curricular y generación de
  artefactos de salida.

  \item \textbf{Diseñar mecanismos de importación deterministas y asistidos por
  modelos de lenguaje.} El objetivo consiste en incorporar a Open CVN currículos ya
  existentes en CVN XML o en CVN PDF, priorizando siempre una vía de extracción
  determinista y reservando el uso de modelos de lenguaje de gran tamaño (del
  inglés, \textit{Large Language Models}, en adelante LLM) como alternativa
  opcional para los casos en que la vía determinista no resulta suficiente. Esta
  vía asistida debe activarse de forma explícita, mantenerse trazable y quedar
  sujeta a validación posterior frente al formato Open CVN antes de aceptarse como
  resultado válido.

  \item \textbf{Verificar el sistema mediante pruebas reproducibles.} La
  evaluación se apoya en pruebas automatizadas, comandos documentados y flujos de
  validación extremo a extremo, con el fin de establecer las garantías reales del
  sistema desarrollado.
\end{enumerate}

La Tabla~\ref{tab:objetivos_capitulos} relaciona estos objetivos con los
capítulos principales en los que se desarrollan.

\begin{table}[htbp]
  \centering
  \footnotesize
  \renewcommand{\arraystretch}{1.55}
  \caption{Relación entre objetivos específicos y capítulos del documento.}
  \label{tab:objetivos_capitulos}
  \begin{tabular}{p{0.12\textwidth} p{0.58\textwidth} p{0.18\textwidth}}
    \hline
    \textbf{Objetivo} & \textbf{Descripción resumida} & \textbf{Capítulos} \\
    \hline
    OE1 & Análisis del ecosistema CVN y sus artefactos oficiales & 2, 3 \\
    OE2 & Estudio de alternativas de representación, modelado y validación & 2 \\
    OE3 & Generación estructural reproducible desde XSD & 4, 5 \\
    OE4 & Normalización de metadatos CVN & 5 \\
    OE5 & Definición de reglas semánticas trazables & 3, 5 \\
    OE6 & Generación de modelos de dominio y artefactos conceptuales & 5 \\
    OE7 & Definición del formato Open CVN JSON & 6 \\
    OE8 & Parser, validador, almacenamiento y exportación & 6 \\
    OE9 & Importación determinista y asistida por LLM & 6 \\
    OE10 & Verificación automatizada y flujos reproducibles & 7 \\
    \hline
  \end{tabular}
\end{table}

\section{Alcance del proyecto}
\label{sec:introduccion_alcance}

El alcance del proyecto se centra en la construcción de una base computacional
para trabajar con información curricular estructurada. El trabajo incluye el
análisis del paquete oficial CVN empleado como fuente, la generación de bindings
estructurales, la normalización de metadatos, la definición de reglas semánticas,
la generación de modelos de dominio, la especificación de un formato JSON propio y
la implementación de herramientas de validación, importación, almacenamiento y
exportación. Estas tareas se han organizado de forma que cada capa conserve
trazabilidad hacia las fuentes originales y pueda verificarse mediante
procedimientos reproducibles.

Dentro del alcance se incluye también la importación de currículos ya existentes
hacia Open CVN a partir de CVN XML y de CVN PDF. Esta importación se resuelve en
primer lugar mediante extracción determinista y solo recurre, de forma opcional y
explícitamente autorizada, a modelos de lenguaje (LLM) cuando la vía determinista
no permite recuperar información suficiente. El resultado de esta vía asistida no
se presenta como una validación semántica completa del currículo importado, sino
como una propuesta que debe superar la validación del formato Open CVN y quedar
disponible para su revisión antes de considerarse definitiva.

Queda fuera del alcance garantizar una conversión completa y sin pérdida de todos
los currículos CVN existentes, ya que dicha garantía requeriría cubrir la totalidad
de variantes, casos particulares y restricciones prácticas del ecosistema oficial.
Tampoco se pretende sustituir la norma CVN ni la infraestructura proporcionada por
FECYT. La propuesta debe entenderse como una capa abierta y reproducible que
facilita el tratamiento de datos curriculares cuando se necesita validación,
transformación, almacenamiento local o exportación a otros formatos.

La herramienta local desarrollada dentro del proyecto cumple una función
demostrativa y operativa. Su finalidad es comprobar que el formato propuesto y las
capas de validación pueden utilizarse en flujos de importación, almacenamiento,
versionado y exportación. No obstante, la contribución principal del TFG no reside
en la interfaz de dicha herramienta, sino en la arquitectura de representación,
normalización y transformación que permite construirla.

\section{Competencias desarrolladas}
\label{sec:introduccion_competencias}

El desarrollo del proyecto se alinea con competencias propias de la
intensificación de Computación. En particular, de acuerdo con la descripción
oficial del trabajo, se abordan las siguientes:

\begin{itemize}
  \item \textbf{[CM1]} Capacidad para tener un conocimiento profundo de los
  principios fundamentales y modelos de la computación y saberlos aplicar para
  interpretar, seleccionar, valorar, modelar y crear nuevos conceptos, teorías,
  usos y desarrollos tecnológicos relacionados con la informática.

  \item \textbf{[CM2]} Capacidad para conocer los fundamentos teóricos de los
  lenguajes de programación y las técnicas de procesamiento léxico, sintáctico y
  semántico asociadas, y saber aplicarlas para la creación, diseño y procesamiento
  de lenguajes.

  \item \textbf{[CM5]} Capacidad para adquirir, obtener, formalizar y representar
  el conocimiento humano en una forma computable para la resolución de problemas
  mediante un sistema informático en cualquier ámbito de aplicación,
  particularmente los relacionados con aspectos de computación, percepción y
  actuación en ambientes o entornos inteligentes.

  \item \textbf{[CM6]} Capacidad para desarrollar y evaluar sistemas interactivos
  y de presentación de información compleja y su aplicación a la resolución de
  problemas de diseño de interacción persona-computadora.
\end{itemize}

Estas competencias se reflejan en el análisis de estructuras oficiales complejas,
la formalización de información curricular en modelos computables, la generación
automática de artefactos, la definición de reglas de validación, la construcción
de una herramienta de gestión y la evaluación técnica del sistema mediante pruebas
reproducibles. Su grado de cumplimiento se retomará en el capítulo de
conclusiones, una vez presentados el desarrollo y los resultados obtenidos.

\section{Contribuciones}
\label{sec:introduccion_contribuciones}

Las principales contribuciones del trabajo pueden resumirse en los siguientes
elementos:

\begin{itemize}
  \item \textbf{Análisis computacional del ecosistema CVN:} se identifican las
  fuentes estructurales, funcionales y auxiliares necesarias para interpretar la
  información curricular representada en CVN, así como las restricciones que
  condicionan su uso automatizado.

  \item \textbf{Arquitectura reproducible por capas:} se separan los bindings
  generados, la normalización de metadatos, la política semántica, los modelos de
  dominio, el modelo conceptual, el formato JSON, la validación y la herramienta
  local, evitando que una única representación concentre todas las
  responsabilidades del sistema.

  \item \textbf{Trazabilidad entre CVN y Open CVN:} se conserva evidencia sobre
  códigos, fuentes, rutas y decisiones semánticas, con el fin de que la
  transformación hacia el nuevo formato mantenga una relación verificable con el
  estándar original.

  \item \textbf{Formato Open CVN JSON:} se define una representación abierta,
  validable y extensible orientada al intercambio, almacenamiento y exportación de
  datos curriculares.

  \item \textbf{Herramientas de validación y uso:} se implementan parsers,
  validadores, almacenamiento local, versionado curricular y exportación a otros
  formatos, lo que permite comprobar la utilidad práctica de la representación
  definida.

  \item \textbf{Importación determinista y asistida por LLM:} se implementan
  vías de importación desde CVN XML y CVN PDF hacia Open CVN, priorizando la
  extracción determinista y reservando el uso de modelos de lenguaje de gran
  tamaño como alternativa opcional, trazable y sujeta a validación posterior para
  los casos en que la extracción determinista resulta insuficiente.

  \item \textbf{Verificación automatizada:} se incorporan pruebas y flujos
  reproducibles para evaluar la consistencia de los artefactos generados y el
  comportamiento de las capas principales del sistema.
\end{itemize}

En conjunto, estas contribuciones no deben entenderse como una sustitución del
ecosistema CVN, sino como una propuesta complementaria que facilita su tratamiento
computacional en contextos donde resultan necesarias la validación formal, la
transformación entre formatos y la gestión automatizada de currículos.

\section{Estructura del documento}
\label{sec:introduccion_estructura}

El resto del documento se organiza de forma progresiva, desde los antecedentes y
el análisis del dominio hasta la implementación, evaluación y discusión de la
solución desarrollada:

\begin{itemize}
  \item \textbf{Capítulo 2: Antecedentes y estado del arte.} Presenta el contexto
  técnico y conceptual necesario para comprender el problema, incluyendo CVN, la
  interoperabilidad curricular, los modelos relacionados, los formatos de
  serialización y las alternativas de modelado.

  \item \textbf{Capítulo 3: Análisis del ecosistema CVN y propuesta de solución.}
  Describe las fuentes oficiales utilizadas, sus limitaciones y la propuesta de
  arquitectura Open CVN derivada de dicho análisis.

  \item \textbf{Capítulo 4: Metodología, herramientas y arquitectura general.}
  Explica el proceso técnico seguido, las herramientas empleadas y la organización
  por capas del sistema.

  \item \textbf{Capítulo 5: Implementación del \textit{pipeline} de generación y
  normalización.} Detalla la generación estructural, la normalización de
  metadatos, la política semántica y la generación de modelos y artefactos
  conceptuales.

  \item \textbf{Capítulo 6: Formato Open CVN, validación y herramienta de
  gestión.} Presenta el formato JSON propuesto, los parsers, validadores, las
  vías de importación determinista y asistida por LLM, los mecanismos de
  almacenamiento local y los flujos de exportación.

  \item \textbf{Capítulo 7: Evaluación, resultados y discusión.} Describe la
  estrategia de verificación, los resultados obtenidos y las garantías reales que
  ofrece el sistema desarrollado.

  \item \textbf{Capítulo 8: Conclusiones, competencias y trabajo futuro.} Resume
  el grado de cumplimiento de los objetivos, relaciona el trabajo con las
  competencias desarrolladas, recoge las limitaciones identificadas y plantea
  líneas futuras.
\end{itemize}
